\chapter{Pemetaan Linear}\label{ch:b1:linmaps}

\omterm{def:b1:logic:map}{Pemetaan} yang pantas ditelaah antara \omterm{def:b1:vspaces:def}{ruang vektor} adalah yang selaras
dengan strukturnya: yaitu \emph{\omterm{def:b1:linmaps:def}{pemetaan linear}}. Adapun kedua
\omterm{def:b1:vspaces:subspace}{subruang} fundamentalnya --- \omterm{def:b1:linmaps:kerim}{kernel} dan \omterm{def:b1:linmaps:kerim}{peta} --- mengukur keinjektifan dan
kesurjektifannya, dan dalam \omterm{def:b1:findim:def}{dimensi hingga} teorema rank--nulitas
mengikat ukurannya menjadi satu hukum kekekalan. Lalu \omterm{def:b1:linmaps:projection}{proyeksi} dan
simetri, kemudian \omterm{def:b1:linmaps:forms}{bentuk linear} dan \omterm{def:b1:linmaps:forms}{hiperbidang}, menutup bab ini.

\section{Definisi dan sifat yang pertama}

\begin{definition}\label{def:b1:linmaps:def}
Misalkan $E, F$ \omterm{def:b1:vspaces:def}{ruang vektor} atas $K$. Sebuah \omterm{def:b1:logic:map}{pemetaan} $u \colon E \to F$ disebut
\emph{linear}\index{pemetaan linear} bila
\[
\forall x, y \in E,\ \forall \lambda \in K, \qquad
u(x + \lambda y) = u(x) + \lambda u(y).
\]
Maka $u(0) = 0$ dan $u(\sum \lambda_i x_i) = \sum \lambda_i u(x_i)$.
Adapun \omterm{def:b1:logic:sets}{himpunan} $\mathcal{L}(E, F)$ berisi \omterm{def:b1:logic:map}{pemetaan} linear sendirinya merupakan \omterm{def:b1:vspaces:def}{ruang vektor};
komposisi \omterm{def:b1:logic:map}{pemetaan} linear bersifat linear, dan bilinear pada pasangannya. Sebuah
\emph{endomorfisma} adalah \omterm{def:b1:logic:map}{pemetaan} linear $E \to E$; sedangkan \emph{isomorfisma} adalah
\omterm{def:b1:logic:map}{pemetaan} linear yang \omterm{def:b1:logic:inj}{bijektif} (yang inversnya lalu otomatis linear);
adapun $u^{-1}$ sebuah isomorfisma, dan komposisi isomorfisma, merupakan
isomorfisma.
\end{definition}

\begin{proof}[Bukti bahwa inversnya linear]
Misalkan $u$ \omterm{def:b1:linmaps:def}{linear} dan \omterm{def:b1:logic:inj}{bijektif}, $y, y' \in F$ dan $\lambda \in K$.
Tetapkanlah $x = u^{-1}(y)$ dan $x' = u^{-1}(y')$. Maka
\[
u\bigl(x + \lambda x'\bigr) = u(x) + \lambda u(x')
= y + \lambda y' ,
\]
lalu menerapkan $u^{-1}$ pada kedua ujungnya: $u^{-1}(y + \lambda y') = x
+ \lambda x' = u^{-1}(y) + \lambda\,u^{-1}(y')$. Jadi tak ada apa pun tentang
$u^{-1}$ yang dihitung: karena kelinearannya terangkut lewat
sifat pendefinisi $u$ belaka --- yaitu pola yang pantas diingat,
sebab struktur sering menumpang bijeksi secara cuma-cuma.
\end{proof}

\begin{proposition}[Sebuah pemetaan linear dikenal pada sebuah basis]\label{prop:b1:linmaps:basis}
Misalkan $(e_1, \dots, e_n)$ sebuah \omterm{def:b1:vspaces:free}{basis} $E$ dan $(v_1, \dots, v_n)$
vektor $F$ yang sebarang. Maka ada tepat satu \omterm{def:b1:linmaps:def}{pemetaan linear} $u \colon
E \to F$ dengan $u(e_i) = v_i$ untuk setiap $i$. Lebih lanjut:
\[
u \text{ injektif} \iff (v_i) \text{ bebas};
\qquad
u \text{ surjektif} \iff (v_i) \text{ membangun } F .
\]
\end{proposition}

\begin{proof}
Keberadaan dan ketunggalannya: setiap $x$ mempunyai \omterm{prop:b1:vspaces:coordinates}{koordinat} tunggal $x = \sum
\lambda_i e_i$ (\cref{prop:b1:vspaces:coordinates}); lalu kelinearannya
memaksa $u(x) = \sum \lambda_i v_i$, dan rumus ini memang mendefinisikan sebuah
\omterm{def:b1:linmaps:def}{pemetaan linear}.

Keinjektifannya: menurut \cref{prop:b1:linmaps:kernel} di bawah, $u$ bersifat
\omterm{def:b1:logic:inj}{injektif} jika dan hanya jika \omterm{def:b1:linmaps:kerim}{kernelnya} sepele. Sekarang $u\bigl(\sum\lambda_i
e_i\bigr) = 0$ persis berarti $\sum\lambda_i v_i = 0$. Jika $(v_i)$
\omterm{def:b1:vspaces:free}{bebas}, maka ini memaksa setiap $\lambda_i = 0$, yakni \omterm{def:b1:linmaps:kerim}{kernelnya}
tereduksi menjadi $0$: jadi \omterm{def:b1:logic:inj}{injektif}. Sedangkan jika $(v_i)$ terikat, sebuah relasi
taksepele $\sum\lambda_i v_i = 0$ menghasilkan vektor taknol
$\sum\lambda_i e_i$ pada \omterm{def:b1:linmaps:kerim}{kernelnya} (karena $(e_i)$ \omterm{def:b1:vspaces:free}{bebas}): jadi tak
\omterm{def:b1:logic:inj}{injektif}. Sehingga kedua syaratnya cocok suku demi suku.

Kesurjektifannya: \omterm{def:b1:linmaps:kerim}{peta} $u$ adalah \omterm{def:b1:logic:sets}{himpunan} semua $\sum\lambda_i
v_i$, yakni tepat $\operatorname{Vect}(v_1, \dots, v_n)$, yang
sama dengan $F$ jika dan hanya jika keluarganya membangun.
\end{proof}

\begin{definition}[Kernel dan peta]\label{def:b1:linmaps:kerim}
Untuk $u \in \mathcal{L}(E, F)$:
\[
\ker u = \{x \in E : u(x) = 0\} \subseteq E,
\qquad
\operatorname{im} u = u(E) \subseteq F ,
\]
keduanya \omterm{def:b1:vspaces:subspace}{subruang} (lewat pemeriksaan langsung dengan
kriterianya).\index{kernel}\index{peta}
\end{definition}

\begin{method}[Kernel dan peta, dalam praktik]
\label{met:b1:linmaps:kerim}
\emph{\omterm{def:b1:linmaps:kerim}{Kernelnya}}: tulislah $u(x) = 0$ sebagai sebuah sistem atas \omterm{prop:b1:vspaces:coordinates}{koordinat}
(atau koefisien) $x$, pecahkanlah, lalu parameterkanlah --- maka \omterm{def:b1:linmaps:kerim}{kernelnya}
keluar dengan sebuah \omterm{def:b1:vspaces:free}{basis} yang menempel
(\cref{met:b1:findim:computedim}). \emph{\omterm{def:b1:linmaps:kerim}{Petanya}}: ia \omterm{def:b1:vspaces:span}{rentang}
\omterm{def:b1:linmaps:kerim}{peta} \emph{sebarang} keluarga pembangun $E$ ---
biasanya sebuah \omterm{def:b1:vspaces:free}{basis}, sehingga $\operatorname{im} u =
\operatorname{Vect}\bigl(u(e_1), \dots, u(e_n)\bigr)$; lalu
hapuslah \omterm{def:b1:linmaps:kerim}{peta} yang berlebihan untuk menyarikan sebuah \omterm{def:b1:vspaces:free}{basis}. \emph{Jalan pintasnya}:
hitunglah yang mana pun di antara keduanya yang lebih mudah lalu peroleh dimensi yang lain
cuma-cuma lewat rank--nulitas
(\cref{thm:b1:linmaps:ranknullity}); dan ketika sebuah calon \omterm{def:b1:vspaces:subspace}{subruang}
yang masuk akal bagi \omterm{def:b1:linmaps:kerim}{petanya} sudah dikenal, membandingkan dimensinya menaikkan
inklusi yang mudah menjadi kesamaan
(\cref{thm:b1:findim:subspaces}). Kedua jalan pintasnya dipakai pada
\cref{ex:b1:linmaps:delta} di bawah.
\end{method}

\begin{proposition}\label{prop:b1:linmaps:kernel}
\omterm{def:b1:logic:map}{Pemetaan} $u$ bersifat \omterm{def:b1:logic:inj}{injektif} $\iff$ $\ker u = \{0\}$; dan $u$ bersifat \omterm{def:b1:logic:inj}{surjektif} $\iff$
$\operatorname{im} u = F$.
\end{proposition}

\begin{proof}
Seperti pada \omterm{def:b1:structures:group}{grup} (\cref{prop:b1:structures:kernel}): $u(x) = u(y) \iff
u(x - y) = 0 \iff x - y \in \ker u$. Adapun butir keduanya merupakan
definisinya.
\end{proof}

\section{Teorema rank--nulitas}

\begin{definition}\label{def:b1:linmaps:rank}
Adapun \emph{rank}\index{rank!pemetaan linear} sebuah $u \in \mathcal{L}(E,
F)$ (dengan $E$ yang \omterm{def:b1:findim:def}{berdimensi hingga}) adalah $\operatorname{rk} u = \dim
\operatorname{im} u$ --- yang juga merupakan rank keluarga
$\bigl(u(e_1), \dots, u(e_n)\bigr)$ bagi sebarang \omterm{def:b1:vspaces:free}{basis} $(e_i)$ pada $E$.
\end{definition}

\begin{theorem}[Rank--nulitas]\label{thm:b1:linmaps:ranknullity}
Misalkan $E$ \omterm{def:b1:findim:def}{berdimensi hingga} dan $u \in \mathcal{L}(E, F)$.
Maka\index{teorema rank--nulitas}
\[
\dim E = \dim \ker u + \operatorname{rk} u .
\]
Lebih tepatnya, jika $S$ sebarang \omterm{def:b1:vspaces:subspace}{subruang} \omterm{def:b1:vspaces:sum}{pelengkap} bagi $\ker u$ di
$E$, maka $u$ terbatasi menjadi \emph{isomorfisma} dari $S$ pada
$\operatorname{im} u$.
\end{theorem}

\begin{proof}
Misalkan $S$ memenuhi $E = \ker u \oplus S$
(\cref{thm:b1:findim:subspaces}), dan misalkan $v \colon S \to
\operatorname{im} u$ pembatasan $u$.

Adapun $v$ bersifat \omterm{def:b1:logic:inj}{injektif}: karena $\ker v = S \cap \ker u = \{0\}$.

Dan $v$ bersifat \omterm{def:b1:logic:inj}{surjektif}: karena sebarang $u(x)$ dengan $x = k + s$ ($k \in \ker u$, $s
\in S$) sama dengan $u(s) = v(s)$.

Jadi $v$ sebuah isomorfisma; dan isomorfisma mengirim \omterm{def:b1:vspaces:free}{basis} ke \omterm{def:b1:vspaces:free}{basis}
(\cref{prop:b1:linmaps:basis}), sehingga $\dim S = \dim\operatorname{im}
u$, lalu $\dim E = \dim\ker u + \dim S$ menyimpulkannya.
\end{proof}

\begin{example}[Membangun sebuah pemetaan menurut spesifikasi]
\label{ex:b1:linmaps:tospec}
Konstruksikanlah $u \in \mathcal{L}(\R^3)$ dengan $\ker u =
\operatorname{Vect}(1,1,1)$ dan $\operatorname{im} u = \{z =
0\}$. Periksa kewarasannya dulu: rank--nulitas menuntut $1 + 2 = 3$ ---
yang selaras, sehingga sebuah penyelesaian mungkin ada. Pilihlah \omterm{def:b1:vspaces:free}{basis} yang disesuaikan dengan
\omterm{def:b1:linmaps:kerim}{kernelnya}, katakanlah $\bigl((1,1,1),\ e_1,\ e_2\bigr)$
(\cref{ex:b1:findim:completion}), lalu tetapkanlah \omterm{def:b1:linmaps:kerim}{petanya}
(\cref{prop:b1:linmaps:basis}):
\[
u(1,1,1) = 0, \qquad u(e_1) = e_1, \qquad u(e_2) = e_2 .
\]
Maka $\ker u \supseteq \operatorname{Vect}(1,1,1)$ dan
$\operatorname{im} u = \operatorname{Vect}(e_1, e_2) = \{z =
0\}$; lalu rank--nulitas memaksa $\dim\ker u = 1$, sehingga \omterm{def:b1:linmaps:kerim}{kernelnya}
tepat garis yang ditetapkan itu. Adapun secara eksplisit, dengan menguraikan $(x, y, z)
= z(1,1,1) + (x - z)e_1 + (y - z)e_2$:
\[
u(x, y, z) = (x - z,\ y - z,\ 0).
\]
Resepnya tersamaratakan: bahwa \omterm{def:b1:linmaps:def}{pemetaan linear} dengan \omterm{def:b1:linmaps:kerim}{kernel} $N$
dan \omterm{def:b1:linmaps:kerim}{peta} $I$ yang ditetapkan ada tepat ketika $\dim N + \dim I = \dim E$
--- dengan keperluannya dari rank--nulitas, dan kecukupannya dari konstruksi
ini.
\end{example}

\begin{corollary}\label{cor:b1:linmaps:samedim}
Jika $\dim E = \dim F$ (yang hingga), maka untuk $u \in \mathcal{L}(E, F)$:
\[
u \text{ injektif} \iff u \text{ surjektif} \iff u \text{
bijektif}.
\]
Khususnya ini berlaku bagi endomorfisma dalam \omterm{def:b1:findim:def}{dimensi hingga}. (Adapun ia
gagal dalam dimensi tak hingga: karena pada $K[X]$, \omterm{def:b1:derivative:def}{turunannya} bersifat \omterm{def:b1:logic:inj}{surjektif}
tetapi tak \omterm{def:b1:logic:inj}{injektif}, sedangkan $P \mapsto XP$ \omterm{def:b1:logic:inj}{injektif} tetapi tak
\omterm{def:b1:logic:inj}{surjektif}.)
\end{corollary}

\begin{proof}
\omterm{def:b1:logic:inj}{Injektif} $\iff \dim\ker u = 0 \iff \operatorname{rk} u = \dim E =
\dim F \iff \operatorname{im} u = F$ (karena \omterm{def:b1:vspaces:subspace}{subruang} yang berdimensi penuh
adalah segalanya, \cref{thm:b1:findim:subspaces}) $\iff$ \omterm{def:b1:logic:inj}{surjektif}.
\end{proof}

\begin{example}[Interpolasi, secara struktural]\label{ex:b1:linmaps:interpolation}
Tetapkanlah $x_0, \dots, x_n$ yang berbeda lalu misalkan $u \colon \R_n[X] \to
\R^{n+1}$, $P \mapsto (P(x_0), \dots, P(x_n))$: yang \omterm{def:b1:linmaps:def}{linear}. Adapun \omterm{def:b1:linmaps:kerim}{kernelnya}
adalah $\{P : \deg P \leq n,\ n+1 \text{ akar}\} = \{0\}$
(\cref{cor:b1:poly:nroots}). Dengan dimensi yang sama $n + 1$: maka $u$ sebuah
isomorfisma --- sehingga keberadaan \emph{dan} ketunggalan interpolan
Lagrange (\cref{thm:b1:poly:lagrange}) dalam satu baris.

Adapun pola satu baris yang sama menangani data yang mencampur nilai dan
\omterm{def:b1:derivative:def}{turunan}: $v \colon \R_3[X] \to \R^4$, $P \mapsto \bigl(P(0),
P'(0), P(1), P'(1)\bigr)$ bersifat \omterm{def:b1:linmaps:def}{linear}, dan \omterm{def:b1:linmaps:kerim}{kernelnya} terdiri atas
\omterm{def:b1:poly:def}{polinomial} berderajat $\leq 3$ dengan akar rangkap di $0$ \emph{dan}
$1$, yakni terbagi oleh $X^2(X-1)^2$ yang berderajat $4$: sehingga hanya $P =
0$. Dengan dimensi yang sama lagi: setiap kuadruple data $(P(0),
P'(0), P(1), P'(1))$ terwujud oleh tepat satu kubik ---
yaitu interpolasi Hermite, yang diberikan oleh perhitungan \omterm{def:b1:linmaps:kerim}{kernel} sebelum
rumus apa pun dituliskan (adapun soal akhir pekan \cref{ch:b1:det}
bertemu determinannya).
\end{example}

\begin{example}[Rank--nulitas dalam kerja: operator selisihnya]
\label{ex:b1:linmaps:delta}
Misalkan $\Delta \colon \R_n[X] \to \R_n[X]$, $P \mapsto P(X+1) -
P(X)$: yang \omterm{def:b1:linmaps:def}{linear}. \omterm{def:b1:linmaps:kerim}{Kernelnya}: jika $\Delta P = 0$, maka $P(0) = P(1) =
P(2) = \dots$, sehingga $P - P(0)$ berakar tak hingga banyak dan
lenyap (\cref{cor:b1:poly:nroots}): jadi $\ker\Delta$ adalah garis
konstantanya. Lalu rank--nulitas: $\operatorname{rk}\Delta = (n + 1) - 1
= n$. Dan karena $\deg \Delta P < \deg P$ untuk $P$ yang takkonstan (sebab suku
puncaknya saling meniadakan), $\operatorname{im}\Delta \subseteq \R_{n-1}[X]$,
yang berdimensi tepat $n$: sehingga inklusinya menjadi kesamaan.
Kesimpulannya, tanpa perhitungan \omterm{def:b1:logic:map}{prapeta} apa pun: bahwa \emph{setiap}
\omterm{def:b1:poly:def}{polinomial} $Q$ berderajat $\leq n - 1$ merupakan selisih $Q = P(X+1)
- P(X)$ --- jadi antiturunan diskretnya ada. (Bandingkanlah
soal akhir pekan \cref{ch:b1:vspaces}, yang di situ $\Delta$
dibalik secara eksplisit pada \omterm{def:b1:vspaces:free}{basis} binomialnya.)
\end{example}

\begin{example}[Pembukuan rank sepanjang sebuah komposisi]
\label{ex:b1:linmaps:rankbookkeeping}
Pada $\R_2[X]$, susunlah \omterm{def:b1:derivative:def}{turunan} $D(P) = P'$ (yang \omterm{def:b1:linmaps:rank}{ber-rank} $2$: dengan
\omterm{def:b1:linmaps:kerim}{peta} $\R_1[X]$, dan \omterm{def:b1:linmaps:kerim}{kernel} konstantanya) dengan dirinya sendiri. Maka $D
\circ D = D^2$ memetakan $P \mapsto P''$, dengan \omterm{def:b1:linmaps:kerim}{peta} $\R_0[X]$: jadi \omterm{def:b1:linmaps:rank}{ber-rank}
$1$. Bandingkanlah dengan batas umumnya: yang kasar memberikan
$\operatorname{rk} D^2 \leq \min(2, 2) = 2$; sedangkan rumus persis
pada \cref{exo:b1:linmaps:12} memperhitungkan kehilangannya dengan tepat,
\[
\operatorname{rk} D^2 = \operatorname{rk} D -
\dim\bigl(\ker D \cap \operatorname{im} D\bigr) = 2 - 1 = 1 ,
\]
karena konstantanya (yaitu \omterm{def:b1:linmaps:kerim}{kernel} $D$ yang luar) duduk di dalam
$\R_1[X]$ (yaitu \omterm{def:b1:linmaps:kerim}{peta} $D$ yang dalam) dengan dimensi $1$. Jadi \omterm{def:b1:linmaps:rank}{rank}
hilang persis di tempat \omterm{def:b1:linmaps:kerim}{kernel} yang luar menyergap \omterm{def:b1:linmaps:kerim}{peta} yang dalam
--- itulah kalimat yang harus diingat ketika \omterm{def:b1:linmaps:rank}{rank} komposisinya berkelakuan buruk.
\end{example}

\begin{example}[Kernel dan peta operator Euler]
\label{ex:b1:linmaps:euler}
Pada $\R_n[X]$, misalkan $u(P) = X\,P'$ (yang \omterm{def:b1:linmaps:def}{linear}: karena pendiferensialan dan
perkalian dengan $X$ demikian). \emph{\omterm{def:b1:linmaps:kerim}{Kernelnya}}: $XP' = 0$ memaksa $P' =
0$ (karena hasil kali \omterm{def:b1:poly:def}{polinomial} lenyap hanya bila salah satu faktornya demikian),
sehingga $\ker u$ adalah garis konstantanya. \emph{\omterm{def:b1:linmaps:kerim}{Petanya}}: pada
\omterm{def:b1:vspaces:free}{basis} monomialnya,
\[
u(X^k) = k\,X^{k} \qquad (k = 0, 1, \dots, n),
\]
sehingga $\operatorname{im} u = \operatorname{Vect}(X, 2X^2, \dots,
nX^n) = \operatorname{Vect}(X, X^2, \dots, X^n)$: yaitu
\omterm{def:b1:poly:def}{polinomial} yang bersuku konstan nol. Periksalah terhadap rank--nulitas:
$\operatorname{rk} u = (n + 1) - 1 = n$, yang memang
dimensi yang ditemukan. Ada dua catatan yang pantas disimpan. Pertama, di sini
$\operatorname{im} u \oplus \ker u = \R_n[X]$ --- tetapi itu
\emph{kebetulan yang membahagiakan} bagi operator ini, bukan sebuah teorema: karena bagi
$v(P) = P'$ yang mirip geseran pada $\R_1[X]$, $\ker v = \operatorname{im}
v = \R_0[X]$ dan jumlahnya tak langsung. Kedua, relasi
$u(X^k) = kX^k$ mengatakan bahwa setiap monomialnya sekadar diskalakan ulang oleh $u$ ---
yaitu \omterm{def:b1:vspaces:free}{basis} yang disesuaikan dengan \omterm{def:b1:logic:map}{pemetaannya}, yakni benih gagasan nilai eigen
yang dikembangkan pada jilid Tahun~2.
\end{example}

\section{Proyeksi dan simetri}

\begin{definition}\label{def:b1:linmaps:projection}
Misalkan $E = F \oplus G$. Adapun \emph{proyeksi pada $F$ sepanjang
$G$}\index{proyeksi} memetakan $x = f + g$ (dengan penguraiannya yang tunggal) ke
$p(x) = f$; sedangkan \emph{simetri}\index{simetri} yang berkaitan adalah $s(x)
= f - g$. Keduanya \omterm{def:b1:linmaps:def}{linear}, dan $s = 2p - \mathrm{id}$.
\end{definition}

\begin{theorem}[Pencirian aljabarnya]\label{thm:b1:linmaps:projchar}
\begin{enumerate}
  \item Sebuah endomorfisma $p$ merupakan \omterm{def:b1:linmaps:projection}{proyeksi} (pada suatu $F$ sepanjang
        suatu $G$) jika dan hanya jika $p \circ p = p$; dan lalu $F =
        \operatorname{im} p = \ker(p - \mathrm{id})$ serta $G = \ker
        p$.
  \item Sebuah endomorfisma $s$ merupakan simetri jika dan hanya jika $s \circ s
        = \mathrm{id}$; dan lalu $E = \ker(s - \mathrm{id}) \oplus
        \ker(s + \mathrm{id})$.
\end{enumerate}
\end{theorem}

\begin{proof}
(1) Sebuah \omterm{def:b1:linmaps:projection}{proyeksi} memenuhi $p(f + g) = f$ dan $p(f) = f$: sehingga $p^2 = p$.
Sebaliknya, misalkan $p^2 = p$; tetapkanlah $F = \operatorname{im} p$, $G = \ker
p$. Maka setiap $x$ tertulis $x = p(x) + (x - p(x))$ dengan $p(x) \in F$ dan
$p\bigl(x - p(x)\bigr) = p(x) - p^2(x) = 0$: sehingga $E = F + G$. Jika $y \in
F \cap G$: maka $y = p(z)$ dan $p(y) = 0$, sehingga $y = p(z) = p^2(z) = p(y) =
0$: jadi jumlahnya langsung, dan $p$ merupakan \omterm{def:b1:linmaps:projection}{proyeksi} pada $F$ sepanjang $G$.
Akhirnya pada $F$: $y = p(z)$ memberikan $p(y) = y$, sehingga $F \subseteq
\ker(p - \mathrm{id})$, dan sebaliknya $p(y) = y$ menaruh $y$ di
\omterm{def:b1:linmaps:kerim}{petanya}.

(2) Korespondensi $s = 2p - \mathrm{id}$, $p = \frac{s +
\mathrm{id}}2$ merupakan bijeksi antara endomorfismanya, dan di bawahnya
\[
s^2 = 4p^2 - 4p + \mathrm{id} = \mathrm{id}
\iff 4p^2 = 4p \iff p^2 = p :
\]
sehingga simetrinya tepat bersesuaian dengan \omterm{def:b1:linmaps:projection}{proyeksinya}. Lalu menerjemahkan
\omterm{def:b1:vspaces:subspace}{subruangnya}: $s(x) = x \iff p(x) = x$, sehingga $\ker(s - \mathrm{id}) =
\operatorname{im} p = F$; sedangkan $s(x) = -x \iff 2p(x) = 0 \iff x
\in \ker p = G$, sehingga $\ker(s + \mathrm{id}) = G$. Jadi \omterm{def:b1:vspaces:sum}{jumlah langsung}
$E = F \oplus G$ pada butir (1) menjadi penguraian yang diumumkan
menjadi vektor yang tetap dan vektor yang terbalik oleh $s$.
\end{proof}

\begin{example}[Sebuah proyeksi dan simetrinya, secara eksplisit]
\label{ex:b1:linmaps:projexplicit}
Di dalam $\R^2$, proyeksikanlah pada $F = \operatorname{Vect}(1,1)$ sepanjang $G =
\operatorname{Vect}(0,1)$. Uraikanlah $(x, y) = a(1,1) + b(0,1)$:
maka \omterm{prop:b1:vspaces:coordinates}{koordinat} pertamanya memberikan $a = x$, dan yang kedua $b = y - x$.
Sehingga
\[
p(x, y) = (x, x),
\qquad
s(x, y) = 2p(x,y) - (x,y) = (x,\ 2x - y).
\]
Periksalah aljabarnya: $p(p(x,y)) = p(x,x) = (x,x)$, dan $s(s(x,y)) =
s(x, 2x - y) = (x, 2x - (2x - y)) = (x, y)$. Secara geometri, $s$
merupakan ``pencerminan miring'' terhadap garis $y = x$ pada arah
tegaknya: karena ia menetapkan $F$ titik demi titik dan membalikkan $G$. Seandainya
kita memproyeksikan pada $F$ yang sama sepanjang $G' =
\operatorname{Vect}(1,-1)$, maka rumusnya akan berubah menjadi
$p'(x,y) = \bigl(\frac{x+y}2, \frac{x+y}2\bigr)$: jadi sebuah \omterm{def:b1:linmaps:projection}{proyeksi}
ditentukan oleh \omterm{def:b1:linmaps:kerim}{petanya} \emph{dan} \omterm{def:b1:linmaps:kerim}{kernelnya}, tak pernah oleh \omterm{def:b1:linmaps:kerim}{petanya}
belaka.
\end{example}

\begin{omfigure}
\begin{tikzpicture}[scale=1.15]
  \draw[->, thin] (-0.4,0) -- (3.4,0) node[below] {$x$};
  \draw[->, thin] (0,-0.4) -- (0,3.8) node[left] {$y$};
  \draw[omDef, very thick] (-0.3,-0.3) -- (3.1,3.1)
    node[right, pos=0.98, font=\small] {$F:\ y = x$};
  \draw[omProp, thick, ->] (2.6,0.2) -- (2.6,0.9)
    node[right, font=\small] {$G$};
  \fill (2,0.5) circle (0.05) node[right, font=\small] {$M$};
  \fill (2,2) circle (0.05) node[left, font=\small] {$p(M)$};
  \fill (2,3.5) circle (0.05) node[left, font=\small] {$s(M)$};
  \draw[omThm, dashed] (2,0.5) -- (2,3.5);
\end{tikzpicture}

{\small \omterm{def:b1:linmaps:projection}{Proyeksi} pada $F = \operatorname{Vect}(1,1)$ sepanjang
$G = \operatorname{Vect}(0,1)$ beserta simetrinya, pada titik $M
= (2,\ 0.5)$: dengan meluncur tegak, $M$ mengenai $F$ di $p(M) = (2,2)$
lalu mendarat di $s(M) = 2p(M) - M = (2,\ 3.5)$, sejauh di atas $F$
(yang diukur sepanjang $G$) sebagaimana $M$ tadi di bawahnya.}
\end{omfigure}

\section{Bentuk linear dan hiperbidang}

\begin{definition}\label{def:b1:linmaps:forms}
Sebuah \emph{bentuk linear}\index{bentuk linear} pada $E$ adalah \omterm{def:b1:linmaps:def}{pemetaan linear}
$\varphi \colon E \to K$. Sedangkan \emph{hiperbidang}\index{hiperbidang}
$E$ (dengan $\dim E = n$) adalah \omterm{def:b1:vspaces:subspace}{subruang} berdimensi $n - 1$.
\end{definition}

\begin{example}[Sebuah bentuk penilaian dan hiperbidangnya]
\label{ex:b1:linmaps:evaluationform}
Pada $\R_2[X]$, penilaian $\varphi(P) = P(2)$ merupakan \omterm{def:b1:linmaps:forms}{bentuk
linear} yang taknol (karena $\varphi(1) = 1$). Adapun \omterm{def:b1:linmaps:kerim}{kernelnya} adalah \omterm{def:b1:linmaps:forms}{hiperbidang}
berisi \omterm{def:b1:poly:def}{polinomial} yang lenyap di $2$, yakni (menurut teorema faktor,
\cref{thm:b1:poly:factor}) kelipatan $X - 2$ di dalam
$\R_2[X]$:
\[
\ker\varphi = \operatorname{Vect}\bigl(X - 2,\ X(X - 2)\bigr),
\qquad \dim = 2 .
\]
Dalam \omterm{prop:b1:vspaces:coordinates}{koordinat} pada $(1, X, X^2)$, $\varphi(a + bX + cX^2) = a +
2b + 4c$: jadi setiap \omterm{def:b1:linmaps:forms}{bentuk linear} pada ruang yang \omterm{def:b1:findim:def}{berdimensi hingga},
begitu sebuah \omterm{def:b1:vspaces:free}{basis} ditetapkan, merupakan ungkapan \omterm{def:b1:linmaps:def}{linear} yang tetap dalam
\omterm{prop:b1:vspaces:coordinates}{koordinatnya} --- sehingga bentuknya menjadi ``vektor baris'', sebagaimana
\cref{ch:b1:matrices} akan membuatnya harfiah, dan baris koefisien
di sini, $(1, 2, 4)$, merupakan baris Vandermonde: jadi bentuk penilaian
adalah cara teori interpolasi pada soal akhir pekan \cref{ch:b1:det}
masuk ke aljabar \omterm{def:b1:linmaps:def}{linear}.
\end{example}

\begin{theorem}\label{thm:b1:linmaps:hyperplanes}
\omterm{def:b1:linmaps:forms}{Hiperbidang} $E$ tepat merupakan \omterm{def:b1:linmaps:kerim}{kernel} \omterm{def:b1:linmaps:forms}{bentuk linear}
yang taknol. Dan dua bentuk taknol berkernel sama jika dan hanya jika keduanya
sebanding.
\end{theorem}

\begin{proof}
Jika $\varphi \neq 0$: maka $\operatorname{rk}\varphi = 1$ (karena \omterm{def:b1:linmaps:kerim}{petanya} \omterm{def:b1:vspaces:subspace}{subruang}
taknol pada $K$), sehingga $\dim\ker\varphi = n - 1$: yaitu sebuah \omterm{def:b1:linmaps:forms}{hiperbidang}.
Sebaliknya, misalkan $H$ sebuah \omterm{def:b1:linmaps:forms}{hiperbidang}, dengan $(e_1, \dots, e_{n-1})$ sebuah \omterm{def:b1:vspaces:free}{basis}
$H$ yang dilengkapi oleh $e_n$: maka bentuk ``\omterm{prop:b1:vspaces:coordinates}{koordinat} terakhir'' berkernel
$H$.

Adapun bentuk yang sebanding berbagi \omterm{def:b1:linmaps:kerim}{kernelnya}. Sebaliknya, andaikan
$\ker\varphi = \ker\psi = H$ lalu pilihlah $a \notin H$: maka setiap $x$ tertulis
$x = h + \lambda a$ (karena $E = H \oplus Ka$), dan
\[
\varphi(x) = \lambda \varphi(a), \qquad \psi(x) = \lambda\psi(a):
\]
sehingga $\varphi = \frac{\varphi(a)}{\psi(a)}\,\psi$.
\end{proof}

\begin{example}\label{ex:b1:linmaps:hyperplane}
Di dalam $K^n$, sebuah \omterm{def:b1:linmaps:forms}{hiperbidang} merupakan \omterm{def:b1:logic:sets}{himpunan} penyelesaian $\{a_1 x_1 + \dots + a_n
x_n = 0\}$ dengan tak semua $a_i$ nol --- yaitu persamaan yang dikenal bagi sebuah
bidang lewat titik asal di $\R^3$. Adapun pada ruang fungsi, bentuk
penilaian $P \mapsto P(1)$ atau $f \mapsto \int_0^1 f$ mendefinisikan \omterm{def:b1:linmaps:forms}{hiperbidang}
$\R_n[X]$, dan $C(\intcc{0}{1})$ (bandingkan \cref{exo:b1:findim:6}).
\end{example}

\begin{example}[Satu hiperbidang, yang dikerjakan tiga cara]
\label{ex:b1:linmaps:hyperplanerun}
Ambillah $\varphi(x, y, z) = x - 2y + 3z$ pada $\R^3$ dan $H =
\ker\varphi$. \emph{Basisnya}: pecahkanlah $x = 2y - 3z$:
\[
(2y - 3z,\ y,\ z) = y\,(2, 1, 0) + z\,(-3, 0, 1),
\]
yaitu dua vektor yang \omterm{def:b1:vspaces:free}{bebas}: sehingga $\dim H = 2$, jadi sebuah \omterm{def:b1:linmaps:forms}{hiperbidang}, sebagaimana
diramalkan \cref{thm:b1:linmaps:hyperplanes} dari $\varphi \neq
0$. \emph{Garis \omterm{def:b1:vspaces:sum}{pelengkapnya}}: sebarang vektor di luar $H$ merentangnya,
misalnya $a = (1, 0, 0)$ (karena $\varphi(a) = 1 \neq 0$); dan
penguraian sebarang $v$ bersifat eksplisit:
\[
v = \underbrace{\bigl(v - \varphi(v)\,a\bigr)}_{\in\,H}
+ \underbrace{\varphi(v)\,a}_{\in\,\operatorname{Vect}(a)},
\]
karena $\varphi\bigl(v - \varphi(v)a\bigr) = \varphi(v) -
\varphi(v)\varphi(a) = 0$. \emph{Kesebandingannya}: jika $\psi(x,y,z)
= -2x + 4y - 6z$, maka $\psi = -2\varphi$ dan keduanya berkernel
$H$; sebaliknya, sebarang bentuk yang lenyap pada $H$ merupakan kelipatan
$\varphi$ (\cref{exo:b1:linmaps:8}) --- jadi persamaan sebuah
\omterm{def:b1:linmaps:forms}{hiperbidang} tunggal sampai skalanya, yaitu fakta yang terus-menerus dipakai bagi
bidang dalam geometri.
\end{example}

\begin{remark}[Jebakan yang lazim]
\label{rem:b1:linmaps:pitfalls}
\emph{\omterm{def:b1:linmaps:kerim}{Kernel} dan \omterm{def:b1:linmaps:kerim}{peta} hidup di ruang yang berbeda}: karena $\ker u
\subseteq E$, $\operatorname{im} u \subseteq F$; sehingga jumlah $\ker u
+ \operatorname{im} u$ hanya bermakna bagi endomorfisma, dan
bahkan lalu ia tak harus langsung (karena $u(x, y) = (y, 0)$ mempunyai $\ker u
= \operatorname{im} u$; adapun \cref{exo:b1:linmaps:7} mencirikan
kapan kelangsungannya berlaku).
\emph{$u^2 = 0$ tak berarti $u = 0$}: karena $u(x,y) =
(y, 0)$ yang sama berkuadrat nol tanpa lenyap --- adapun yang sungguh dikatakan $u^2 = 0$
adalah $\operatorname{im} u \subseteq \ker u$
(\cref{exo:b1:linmaps:5}).
\emph{\omterm{def:b1:logic:inj}{Injektif} $\iff$ \omterm{def:b1:logic:inj}{surjektif} menuntut \omterm{def:b1:findim:def}{dimensi hingga} yang
sama}: karena pada $K[X]$, \omterm{def:b1:derivative:def}{turunannya} bersifat \omterm{def:b1:logic:inj}{surjektif} dan tak
\omterm{def:b1:logic:inj}{injektif}, sedangkan $P \mapsto XP$ \omterm{def:b1:logic:inj}{injektif} dan tak \omterm{def:b1:logic:inj}{surjektif}
(\cref{cor:b1:linmaps:samedim}); dan di antara ruang yang
berdimensi \emph{berbeda}, satu implikasinya sekadar
mustahil (karena $\operatorname{rk} u \leq \min(\dim E, \dim F)$).
\emph{Menetapkan \omterm{def:b1:linmaps:kerim}{peta} berjalan pada sebuah \omterm{def:b1:vspaces:free}{basis}, bukan pada keluarga apa pun}:
karena menuntut $u(1, 0) = a$, $u(0, 1) = b$, $u(1, 1) = c$
menentukan $u$ secara berlebihan kecuali bila $c = a + b$; jadi \omterm{def:b1:linmaps:def}{pemetaan linear} \omterm{def:b1:vspaces:free}{bebas} pada sebuah
\omterm{def:b1:vspaces:free}{basis}, tetapi terbelenggu di mana pun selainnya.
\emph{\omterm{def:b1:linmaps:rank}{Rank} tak terawetkan oleh komposisi}: karena ia hanya dapat turun,
$\operatorname{rk}(vu) \leq \min(\operatorname{rk} u,
\operatorname{rk} v)$ (\cref{exo:b1:linmaps:4}), dengan kehilangannya yang persis
terukur pada \cref{exo:b1:linmaps:12}.
\end{remark}

\begin{remark}[Ke mana pemetaan ini pergi]
\label{rem:b1:linmaps:whereused}
\omterm{def:b1:linmaps:def}{Pemetaan linear} segera akan menjadi \emph{matriks}: karena begitu basisnya
ditetapkan, \cref{ch:b1:matrices} menyandikan setiap $u \in \mathcal{L}(E,
F)$ lewat sebuah larik persegi panjang, dan komposisinya menjadi hasil kali
matriks --- lalu rank--nulitas menggerakkan teori sistem
\omterm{def:b1:linmaps:def}{linear} pada \cref{ch:b1:det}. Adapun \omterm{def:b1:linmaps:projection}{proyeksi} kembali pada
\cref{ch:b1:euclid} dalam kasus khususnya yang paling berguna, yaitu
\omterm{def:b1:linmaps:projection}{proyeksi} \emph{ortogonal}, yang di situ \omterm{def:b1:linmaps:kerim}{kernelnya} dipilih
tegak lurus terhadap \omterm{def:b1:linmaps:kerim}{petanya}. Sedangkan soal akhir pekan di bawah mendorong
aljabar proyektornya sejauh yang dijangkau perkakas tahun pertama, sampai
lema Fitting; adapun jilid Tahun~2 melangkah lebih jauh dengan trace
dan dengan teori nilai eigen, yang bagi teori itu proyektor pada \omterm{def:b1:vspaces:subspace}{subruang}
yang stabil merupakan balok bangunan dasarnya.
\end{remark}

\begin{remark}[Cakrawala di dalam Buku 3: rank--nulitas tiga kali lagi]
\label{rem:b1:linmaps:perspectives}
Hukum kekekalan $\dim E = \dim\ker u + \operatorname{rk} u$
akan dibaca ulang tiga kali sebelum jilid ini berakhir. Pada
\cref{ch:b1:det} ia menjadi bentuk \omterm{def:b1:logic:sets}{himpunan} penyelesaiannya: bahwa
sistem yang selaras dengan $p$ anu dan \omterm{def:b1:linmaps:rank}{ber-rank} $r$ mempunyai \omterm{def:b1:logic:sets}{himpunan}
penyelesaian berdimensi $p - r$ --- yaitu dimensi \omterm{def:b1:linmaps:kerim}{kernel} yang menyamar. Pada
\cref{ch:b1:euclid} ia membelah secara ortogonal, $\dim F + \dim
F^\perp = \dim E$, lalu menggerakkan setiap perhitungan jarak. Sedangkan pada
soal akhir pekan \cref{ch:b1:multivar}, ia menjadi
akuntan kuadrat terkecilnya: yaitu $n$ pengamatan, $2$ parameter yang
dicocokkan, dan $n - 2$ dimensi sisanya, dengan kesamaan Pythagoras
$\norm b^2 = \norm p^2 + \norm{b - p}^2$ sebagai
bayangan Euclid rank--nulitas. Jadi satu teorema, empat kostum.
\end{remark}

\section{Latihan}

\begin{exercise}[$\star$]\label{exo:b1:linmaps:1}
\omterm{def:b1:logic:map}{Pemetaan} mana yang \omterm{def:b1:linmaps:def}{linear}?
\begin{enumerate}
  \item $\R^2 \to \R^2$, $(x, y) \mapsto (x + y, x - 2y)$;
  \item $\R^2 \to \R$, $(x, y) \mapsto xy$;
  \item $\R[X] \to \R[X]$, $P \mapsto P' + XP$;
  \item $\mathcal{F}(\R,\R) \to \R$, $f \mapsto f(3)$.
\end{enumerate}
\end{exercise}

\begin{exercise}[$\star$]\label{exo:b1:linmaps:2}
Misalkan $u \colon \R^3 \to \R^3$, $(x,y,z) \mapsto (x + y - z,\; 2x + y
+ z,\; 3x + 2y)$. Tentukanlah $\ker u$ (yaitu \omterm{def:b1:vspaces:free}{basis} dan dimensinya),
$\operatorname{rk} u$, serta sebuah \omterm{def:b1:vspaces:free}{basis} $\operatorname{im} u$. Apakah $u$
\omterm{def:b1:logic:inj}{injektif}? Dan \omterm{def:b1:logic:inj}{surjektif}?
\end{exercise}

\begin{exercise}[$\star$]\label{exo:b1:linmaps:3}
Misalkan $u \colon \R_n[X] \to \R_n[X]$, $P \mapsto P - P'$. Buktikan bahwa
$u$ sebuah isomorfisma: sekali lewat $\ker u$, sekali lagi dengan menunjukkan
inversnya \emph{(tinjaulah $P + P' + P'' + \dots$)}.
\end{exercise}

\begin{exercise}[$\star$]\label{exo:b1:linmaps:4}
Misalkan $u \in \mathcal{L}(E, F)$ dan $v \in \mathcal{L}(F, G)$, dengan ruang yang
\omterm{def:b1:findim:def}{berdimensi hingga}. Buktikan:
\[
\operatorname{rk}(v \circ u) \leq
\min\bigl(\operatorname{rk} u,\ \operatorname{rk} v\bigr).
\]
\end{exercise}

\begin{exercise}[$\star\star$]\label{exo:b1:linmaps:5}
Misalkan $u$ sebuah endomorfisma $E$ (yang \omterm{def:b1:findim:def}{berdimensi hingga}) dengan $u^2 =
0$. Buktikan bahwa $\operatorname{im} u \subseteq \ker u$, sehingga
$\operatorname{rk} u \leq \frac{\dim E}{2}$. Untuk $E = \R^2$, berikanlah sebuah
contoh dengan kesamaannya.
\end{exercise}

\begin{exercise}[$\star\star$]\label{exo:b1:linmaps:6}
Misalkan $p, q$ \omterm{def:b1:linmaps:projection}{proyeksi} $E$ dengan $p \circ q = q \circ p$.
Buktikan bahwa $p \circ q$ sebuah \omterm{def:b1:linmaps:projection}{proyeksi}, dengan
\[
\operatorname{im}(pq) = \operatorname{im} p \cap \operatorname{im}
q ,
\qquad
\ker (pq) = \ker p + \ker q .
\]
\end{exercise}

\begin{exercise}[$\star\star$]\label{exo:b1:linmaps:7}
Misalkan $u \in \mathcal{L}(E)$, dengan $E$ yang \omterm{def:b1:findim:def}{berdimensi hingga}. Buktikan
kesetaraan antara:
\begin{enumerate}
  \item $E = \ker u \oplus \operatorname{im} u$;
  \item $\ker u = \ker u^2$;
  \item $\operatorname{im} u = \operatorname{im} u^2$.
\end{enumerate}
\end{exercise}

\begin{exercise}[$\star\star$]\label{exo:b1:linmaps:8}
Misalkan $\varphi, \psi$ \omterm{def:b1:linmaps:forms}{bentuk linear} pada $E$ dengan $\ker\varphi
\subseteq \ker\psi$. Buktikan $\psi = \lambda\varphi$ untuk suatu
$\lambda \in K$ (termasuk kasus yang merosot).
\end{exercise}

\begin{exercise}[$\star\star\star$]\label{exo:b1:linmaps:9}
Misalkan $u \in \mathcal{L}(E)$ dengan $\dim E = n$, dan andaikan $u^n = 0$
tetapi $u^{n-1} \neq 0$ (yaitu endomorfisma yang \emph{nilpoten maksimal}).
Pilihlah $x$ dengan $u^{n-1}(x) \neq 0$; lalu buktikan bahwa $\bigl(x, u(x), \dots,
u^{n-1}(x)\bigr)$ merupakan \omterm{def:b1:vspaces:free}{basis} $E$. \emph{(Terapkanlah pangkat $u$ pada sebuah
kombinasi nol, dimulai dengan $u^{n-1}$.)}
\end{exercise}

\begin{exercise}[$\star\star\star$]\label{exo:b1:linmaps:10}
Misalkan $f \in \mathcal{L}(\R^n)$ dengan $f \circ f = -\mathrm{id}$.
\begin{enumerate}
  \item Buktikan bahwa $f$ sebuah isomorfisma dan bahwa tak ada $x \neq 0$
        yang memenuhi $f(x) = \lambda x$ dengan $\lambda \in \R$.
  \item Buktikan bahwa $n$ genap. \emph{Petunjuk: pilihlah $x_1 \neq 0$; tunjukkan
        bahwa $\operatorname{Vect}(x_1, f(x_1))$ merupakan bidang yang stabil di bawah
        $f$; lalu pilihlah $x_2$ di luarnya dan ulangi, sambil membuktikan bahwa
        $\bigl(x_1, f(x_1), x_2, f(x_2), \dots\bigr)$ tetap \omterm{def:b1:vspaces:free}{bebas}.}
\end{enumerate}
\end{exercise}

\begin{exercise}[$\star\star$]\label{exo:b1:linmaps:11}
Misalkan $u, v \in \mathcal{L}(E, F)$, dengan ruang yang \omterm{def:b1:findim:def}{berdimensi hingga}.
Buktikan batas dua sisinya
\[
\abs{\operatorname{rk} u - \operatorname{rk} v}
\;\leq\; \operatorname{rk}(u + v)
\;\leq\; \operatorname{rk} u + \operatorname{rk} v .
\]
\emph{(Untuk batas atasnya, bandingkanlah $\operatorname{im}(u+v)$ dengan
$\operatorname{im} u + \operatorname{im} v$; sedangkan untuk yang bawah,
terapkanlah batas atasnya dengan cerdik.)}
\end{exercise}

\begin{exercise}[$\star\star\star$]\label{exo:b1:linmaps:12}
(Ketaksamaan Frobenius) Misalkan $u \in \mathcal{L}(E, F)$, $w \in
\mathcal{L}(F, G)$ dan $v \in \mathcal{L}(G, H)$, dengan semua ruangnya
\omterm{def:b1:findim:def}{berdimensi hingga}. Buktikan rumus yang persis
\[
\operatorname{rk}(v \circ w) = \operatorname{rk} w -
\dim\bigl(\ker v \cap \operatorname{im} w\bigr),
\]
lalu simpulkanlah ketaksamaan Frobenius
\[
\operatorname{rk}(v \circ w) + \operatorname{rk}(w \circ u)
\;\leq\; \operatorname{rk} w + \operatorname{rk}(v \circ w \circ
u) .
\]
Periksalah bahwa kasus $w = \mathrm{id}_F$ merupakan ketaksamaan
Sylvester, yang dibuktikan dalam bentuk matriks pada \cref{exo:b1:matrices:10}.
\end{exercise}

\section{Soal: kalkulus proyektor dan lema Fitting}

\begin{problem}\label{pb:b1:linmaps:1}
\omterm{def:b1:linmaps:projection}{Proyeksi} adalah endomorfisma yang dihasilkan \omterm{def:b1:vspaces:sum}{jumlah langsung}, dan
sebaliknya: karena setiap kesamaan $E = F_1 \oplus \dots \oplus F_k$
diam-diam merupakan sebuah keluarga proyektor yang berjumlah identitas. Soal
ini mengembangkan kamus itu --- yaitu aljabar satu
proyektor, dua proyektor, dan $k$ proyektor --- lalu menerapkan gagasan
kestabilan yang sama pada sebarang endomorfisma dan membuktikan
\emph{lema Fitting}: bahwa setiap endomorfisma sebuah ruang yang
\omterm{def:b1:findim:def}{berdimensi hingga} terbelah menjadi bagian nilpoten dan
bagian yang dapat dibalik. Di sepanjang soal ini, $E$ adalah \omterm{def:b1:vspaces:def}{ruang vektor} atas $K$
yang berdimensi $n$, dan \emph{proyektor} berarti $p \in \mathcal{L}(E)$
dengan $p^2 = p$ (\cref{thm:b1:linmaps:projchar}).
\index{proyeksi}\index{lema Fitting}\index{nilpoten}

\medskip
\noindent\textbf{Bagian I --- Aljabar di sekitar satu proyektor.}
Misalkan $p$ sebuah proyektor, dengan $p \neq 0$, $p \neq \mathrm{id}$.
\begin{enumerate}
  \item Tunjukkan bahwa $\mathrm{id} - p$ sebuah proyektor lalu kenalilah
        $\operatorname{im}(\mathrm{id} - p)$ dan
        $\ker(\mathrm{id} - p)$.
  \item Hitunglah $(\lambda\,\mathrm{id} + \mu\,p)^2$ lalu tentukanlah
        semua pasangan $(\lambda, \mu) \in K^2$ yang untuknya
        $\lambda\,\mathrm{id} + \mu\,p$ merupakan proyektor.
  \item Tunjukkan bahwa bidang $\operatorname{Vect}(\mathrm{id}, p)$
        pada $\mathcal{L}(E)$ stabil terhadap komposisi, dan bahwa
        untuk setiap \omterm{def:b1:poly:def}{polinomial} $Q \in K[X]$,
        \[
        Q(p) = Q(0)\,\mathrm{id} + \bigl(Q(1) -
        Q(0)\bigr)\,p .
        \]
  \item Tentukanlah untuk $(\lambda, \mu)$ yang mana \omterm{def:b1:logic:map}{pemetaan}
        $\lambda\,\mathrm{id} + \mu\,p$ dapat dibalik, lalu berikanlah
        inversnya dalam bentuk $\alpha\,\mathrm{id} + \beta\,p$.
        Tafsirkanlah jawabannya lewat tindakan
        $\lambda\,\mathrm{id} + \mu\,p$ pada $\operatorname{im} p$
        dan pada $\ker p$.
  \item Misalkan $p'$ proyektor lain dengan \emph{\omterm{def:b1:linmaps:kerim}{peta} yang sama}
        $\operatorname{im} p' = \operatorname{im} p$. Tunjukkan bahwa
        $p\,p' = p'$ dan $p'\,p = p$. Apa yang dikatakan kesamaan ini
        tentang menyusun \omterm{def:b1:linmaps:projection}{proyeksi} pada \omterm{def:b1:vspaces:subspace}{subruang} yang sama
        sepanjang \omterm{def:b1:linmaps:kerim}{kernel} yang berbeda?
\end{enumerate}

\medskip
\noindent\textbf{Bagian II --- Dua proyektor.} Misalkan $p, q$
proyektor $E$; dan anggaplah karakteristiknya bukan $2$ (yang benar
bagi $K = \R, \C$).
\begin{enumerate}[resume]
  \item Andaikan $p + q$ sebuah proyektor. Dengan menjabarkan $(p + q)^2$,
        tunjukkan $pq + qp = 0$; lalu dengan menyusunnya bersama $p$ di kiri, kemudian
        di kanan, simpulkanlah $pq = qp$, lalu tutuplah dengan $pq = qp =
        0$.
  \item Sebaliknya, andaikan $pq = qp = 0$. Tunjukkan bahwa $p + q$ sebuah
        proyektor, dengan
        \[
        \operatorname{im}(p + q) = \operatorname{im} p \oplus
        \operatorname{im} q,
        \qquad
        \ker(p + q) = \ker p \cap \ker q .
        \]
  \item Tunjukkan bahwa $p - q$ merupakan proyektor jika dan hanya jika $pq = qp
        = q$. \emph{(Terapkanlah pertanyaan 6--7 pada $\mathrm{id} - p$
        dan $q$.)}
  \item Tunjukkan makna geometri $pq = qp = q$: bahwa ia berlaku jika
        dan hanya jika $\operatorname{im} q \subseteq
        \operatorname{im} p$ dan $\ker p \subseteq \ker q$. (Lalu
        kita menulis $q \leq p$: ``$q$ memproyeksikan pada yang lebih sedikit, sepanjang
        yang lebih banyak''.)
  \item Sekarang misalkan $p, q$ berkomutasi. Ingat kembali dari
        \cref{exo:b1:linmaps:6} bahwa $pq$ merupakan proyektor pada
        $\operatorname{im} p \cap \operatorname{im} q$ sepanjang
        $\ker p + \ker q$. Tunjukkan bahwa $r = p + q - pq$ sebuah
        proyektor dengan
        \[
        \operatorname{im} r = \operatorname{im} p +
        \operatorname{im} q,
        \qquad
        \ker r = \ker p \cap \ker q .
        \]
        \emph{(Tinjaulah $\mathrm{id} - r = (\mathrm{id} -
        p)(\mathrm{id} - q)$.)}
\end{enumerate}

\medskip
\noindent\textbf{Bagian III --- Penguraian identitasnya.}
\begin{enumerate}[resume]
  \item Misalkan $E = F_1 \oplus \dots \oplus F_k$ dan, untuk $x = x_1 +
        \dots + x_k$ (dengan penguraiannya yang tunggal, $x_i \in F_i$), tetapkanlah
        $p_i(x) = x_i$. Tunjukkan bahwa setiap $p_i$ sebuah proyektor, bahwa
        $p_i p_j = 0$ untuk $i \neq j$, dan bahwa $p_1 + \dots + p_k
        = \mathrm{id}$; lalu kenalilah $\operatorname{im} p_i$ dan
        $\ker p_i$.
  \item Sebaliknya, misalkan $p_1, \dots, p_k \in \mathcal{L}(E)$
        memenuhi $p_1 + \dots + p_k = \mathrm{id}$ dan $p_i p_j =
        0$ untuk setiap $i \neq j$. Tunjukkan bahwa setiap $p_i$ sebuah
        proyektor dan bahwa $E = \operatorname{im} p_1 \oplus
        \dots \oplus \operatorname{im} p_k$.
  \item Untuk dua proyektor dengan $p + q = \mathrm{id}$: tunjukkan bahwa $pq
        = qp = 0$ berlaku secara otomatis.
  \item Untuk tiga proyektor dengan $p + q + r = \mathrm{id}$: tunjukkan
        bahwa $p + q$ sebuah proyektor, lalu simpulkanlah dari pertanyaan 6
        bahwa \emph{semua} hasil kali berpasangannya lenyap --- sehingga $E =
        \operatorname{im} p \oplus \operatorname{im} q \oplus
        \operatorname{im} r$, tanpa hipotesis apa pun atas hasil kalinya.
  \item Untuk $k$ proyektor dengan $p_1 + \dots + p_k =
        \mathrm{id}$: tunjukkan lebih dulu bahwa untuk sebarang \omterm{def:b1:vspaces:subspace}{subruang}, $\dim
        (F_1 + \dots + F_k) \leq \dim F_1 + \dots + \dim F_k$,
        dengan kesamaannya jika dan hanya jika jumlahnya langsung; lalu
        tunjukkan $E = \operatorname{im} p_1 + \dots +
        \operatorname{im} p_k$, lalu buktikanlah bahwa \emph{jika}
        lebih lanjut $\sum_i \operatorname{rk} p_i \leq n$, maka jumlahnya
        langsung dan $p_i p_j = 0$ untuk $i \neq j$.
\end{enumerate}

\medskip
\noindent\textbf{Bagian IV --- \omterm{def:b1:linmaps:kerim}{Kernel} yang berulang: lema Fitting.}
Misalkan $u \in \mathcal{L}(E)$, dengan $\dim E = n$.
\begin{enumerate}[resume]
  \item Tunjukkan kedua rantainya, yang sah untuk setiap $k \geq 0$:
        \[
        \ker u^k \subseteq \ker u^{k+1},
        \qquad
        \operatorname{im} u^{k+1} \subseteq \operatorname{im}
        u^k .
        \]
  \item Tunjukkan bahwa jika $\ker u^{r} = \ker u^{r+1}$ untuk suatu $r$,
        maka $\ker u^{k} = \ker u^{r}$ untuk setiap $k \geq r$; lalu nyatakanlah
        dan buktikanlah kestabilan yang serupa bagi \omterm{def:b1:linmaps:kerim}{petanya}.
  \item Simpulkan bahwa ada bilangan bulat terkecil $r$ dengan $\ker
        u^{r} = \ker u^{r+1}$, bahwa $r \leq n$, dan bahwa \omterm{def:b1:linmaps:kerim}{petanya}
        stabil pada $r$ yang sama.
  \item (Lema Fitting) Buktikan bahwa
        \[
        E \;=\; \ker u^{r} \,\oplus\, \operatorname{im} u^{r} .
        \]
  \item Tunjukkan bahwa kedua \omterm{def:b1:vspaces:subspace}{subruangnya} stabil di bawah $u$, bahwa
        pembatasan $u$ pada $\ker u^{r}$ bersifat nilpoten, dan bahwa
        pembatasan $u$ pada $\operatorname{im} u^{r}$ merupakan
        isomorfisma $\operatorname{im} u^{r}$: sehingga setiap
        endomorfisma, pada sebuah \omterm{def:b1:vspaces:sum}{jumlah langsung} yang kanonik, bersifat ``nilpoten
        tambah dapat dibalik''.
  \item Misalkan $\pi$ proyektor pada $\ker u^{r}$ sepanjang
        $\operatorname{im} u^{r}$. Tunjukkan bahwa $\pi \circ u = u
        \circ \pi$.
\end{enumerate}

\medskip
\noindent\textbf{Bagian V --- Sebuah kasus yang dikerjakan, dan sintesisnya.}
\begin{enumerate}[resume]
  \item Misalkan $u(x, y, z) = (y, 0, z)$ pada $\R^3$. Hitunglah $u^2$ dan
        $u^3$, tentukanlah indeks kestabilannya $r$, kedua
        \omterm{def:b1:vspaces:subspace}{subruang} $\ker u^{r}$ dan $\operatorname{im} u^{r}$,
        proyektor Fittingnya $\pi$, lalu periksalah pada rumusnya bahwa
        $\pi u = u\pi$ dan bahwa $u$ bersifat nilpoten pada satu faktornya,
        dan \omterm{def:b1:logic:inj}{bijektif} pada yang lain.
  \item Tunjukkan kesetaraannya: bahwa $u$ nilpoten $\iff$ $\ker u^{r} =
        E$ $\iff$ $\pi = \mathrm{id}$; lalu simpulkanlah bahwa
        endomorfisma nilpoten sebuah ruang berdimensi $n$
        selalu memenuhi $u^{n} = 0$ (sehingga indeks kenilpotenannya tak pernah
        melampaui dimensinya).
  \item (Ketunggalan) Andaikan $E = A \oplus B$ dengan $A, B$ yang stabil
        di bawah $u$, dengan pembatasan $u|_A$ yang nilpoten dan $u|_B$ yang
        \omterm{def:b1:logic:inj}{bijektif}. Buktikan $A = \ker u^{r}$ dan $B =
        \operatorname{im} u^{r}$: sehingga penguraian Fittingnya
        tunggal.
  \item Sintesis, dalam empat kalimat: kamus apa yang ditegakkan Bagian III
        antara \omterm{def:b1:vspaces:sum}{jumlah langsung} dan keluarga
        proyektor; mengapakah pertanyaan 14 tak memerlukan hipotesis atas hasil kalinya
        sedangkan pertanyaan 15 memerlukan hipotesis \omterm{def:b1:linmaps:rank}{rank}
        (dan perkakas Tahun~2 yang mana, yaitu trace, yang menghapusnya); dalam
        arti apa lema Fitting merupakan versi yang distabilkan dari
        \cref{exo:b1:linmaps:7}; dan menjadi apa kedua faktor Fittingnya
        pada teori nilai eigen jilid Tahun~2
        itu. Namailah teorema yang dibuktikan pada Bagian IV.

\end{enumerate}
\end{problem}
